\documentclass{report}

\begin{document}
\title{Labwork 2: Get to know your GPU}
\author{Philippe Olivier Schriqui-Duque \\ ID: 2640042}
\date{}
\maketitle

\section*{Introduction}

The goal of this labwork is to get to know my GPU with Numba: its name, the number of
multiprocessors, the number of cores and the memory size.
Before writing CUDA programs it is important to know the hardware we use, because the number
of cores and the memory limit how much work we can give to the GPU.
I work on my laptop with Kali Linux in WSL2, which has an NVIDIA GPU, and I use Python with the
\texttt{numba} library.

\section*{Code}

I use the \texttt{cuda} module of Numba to get the information of my GPU.
First I call \texttt{cuda.detect()} to list the GPUs that Numba can see. I only have one,
so I select it with its id 0 and get the device object. This object has the name and the id
of the GPU:

\begin{verbatim}
from numba import cuda

cuda.detect()
cuda.select_device(0)
device = cuda.get_current_device()
print(f"Device name: {device.name}")
print(f"Device id: {device.id}")
\end{verbatim}

A GPU is made of several streaming multiprocessors (SM), and each SM has a group of CUDA cores.
The number of SM is an attribute of the device:

\begin{verbatim}
MPC = device.MULTIPROCESSOR_COUNT
cores_per_sm = 128
total_cores = MPC * cores_per_sm
\end{verbatim}

Numba does not give the number of cores, so I multiply the number of SM by the cores per SM.
I put 128 directly because I know my GPU: it is an RTX 3070 Laptop (compute capability 8.6),
which has 128 cores per SM. On another GPU this number could be different.

\[
\mbox{cores} = \mbox{SM} \times \mbox{cores per SM} = 40 \times 128 = 5120
\]

The memory size is not an attribute of the device. I get it from the current CUDA context,
which gives the free and the total memory in bytes. I divide by $1024^3$ to have it in GB:

\begin{verbatim}
free, total = cuda.current_context().get_memory_info()
print(f"Memory size: {total / 1024**3:.2f} GB")
\end{verbatim}

\[
\mbox{memory} = \frac{8\,589\,410\,304 \mbox{ bytes}}{1024^3} \approx 8 \mbox{ GB}
\]

I also printed some extra information with other attributes of the device:
\texttt{CLOCK\_RATE}, \texttt{MAX\_THREADS\_PER\_BLOCK}, \texttt{WARP\_SIZE} and
\texttt{MEMORY\_CLOCK\_RATE}. The clock rates are given in kHz, so I divide them by 1000 to get MHz.

\section*{Output}

This is the result when I run \texttt{python labwork2.py}:

\begin{verbatim}
Device name: NVIDIA GeForce RTX 3070 Laptop GPU
Device id: 0
Multiprocessor count: 40
Total cores: 5120
Memory size: 8.00 GB
Clock rate: 1290.00 MHz
Max threads per block: 1024
Warp size: 32
Memory clock: 6001.00 MHz
\end{verbatim}

The results are the same as the official specifications of the RTX 3070 Laptop.
The max threads per block (1024) and the warp size (32) will be useful later, when we choose
the block size of our kernels.

\section*{Conclusion}

In this labwork I learned how to use Numba to get information about my GPU.
My GPU is an RTX 3070 Laptop with 40 SM, 5120 cores and 8 GB of memory.
It has a lot more cores than my CPU (16 threads), so it will be useful for the next labworks
where we run the same simple operation on many pixels at the same time.

\end{document}
